\DocumentMetadata{
  pdfversion=2.0,
  pdfstandard=ua-2,
  lang=en-US,
  tagging=on,
  tagging-setup={math/setup=mathml-SE,tabsorder=structure}
}
\documentclass[11pt,letterpaper]{article}
\usepackage[margin=0.68in,includefoot]{geometry}
\usepackage{newcomputermodern}
\usepackage{amsmath,amscd,mathtools}
\usepackage{tikz,adjustbox}
\providecommand{\square}{\mdlgwhtsquare}
\usepackage[hidelinks]{hyperref}
\hypersetup{
  pdftitle={Functional Analysis PhD Exam, January 2011},
  pdfauthor={University of Florida Department of Mathematics},
  pdfsubject={Graduate mathematics examination},
  pdfdisplaydoctitle=true,
  bookmarksopen=true
}
\setcounter{secnumdepth}{0}
\setlength{\parindent}{0pt}
\setlength{\parskip}{0.28em}
\setlength{\emergencystretch}{3em}
\setlength{\leftmargini}{2.15em}
\setlength{\leftmarginii}{2.65em}
\setlength{\leftmarginiii}{3em}
\renewcommand{\labelenumi}{\textbf{\theenumi.}}
\pagestyle{empty}
\raggedbottom
\allowdisplaybreaks
\newenvironment{examproblems}[1]{%
  \begin{enumerate}%
  \setcounter{enumi}{#1}%
  \setlength{\topsep}{0.35em}%
  \setlength{\partopsep}{0pt}%
  \setlength{\itemsep}{0.48em}%
  \setlength{\parsep}{0.18em}%
}{\end{enumerate}}
\newenvironment{examsubparts}{%
  \begin{enumerate}%
  \setlength{\topsep}{0.18em}%
  \setlength{\partopsep}{0pt}%
  \setlength{\itemsep}{0.16em}%
  \setlength{\parsep}{0.1em}%
}{\end{enumerate}}
\newenvironment{examinstructions}{%
  \par\begingroup\itshape
}{%
  \par\endgroup\vspace{0.18em}
}
\makeatletter
\renewcommand\section{\@startsection{section}{1}{0pt}%
  {-1.25ex plus -.25ex minus -.1ex}%
  {0.55ex plus .1ex}%
  {\normalfont\large\bfseries}}
\renewcommand{\@maketitle}{%
  \newpage\null\vskip -1em%
  {\footnotesize\bfseries
    \href{https://www.ufl.edu/}{University of Florida}, \href{https://math.ufl.edu/}{Department of Mathematics}\par}%
  \vskip -0.5em%
  \tagstructbegin{tag=Title}%
    {\LARGE\bfseries\href{https://gma.math.ufl.edu/exams/functional-analysis-phd/}{Functional Analysis PhD Exam}, January 2011\par}%
  \tagstructend%
  \vskip 0.25em%
  \tagmcbegin{artifact}%
    \hrule height 1pt%
  \tagmcend%
  \vskip 1em%
}
\makeatother
\title{Functional Analysis PhD Exam, January 2011}
\author{}
\date{}
\begin{document}
\maketitle
\thispagestyle{empty}
\begin{examinstructions}
Do only 5 of the 8 problems.
\end{examinstructions}
\begin{examproblems}{0}
\item
State the Krein–Milman Theorem and sketch the proof.
\item
Prove that if~\(A\) is an element of a Banach algebra \(\mathcal U\), then the spectrum of~\(A\) is a nonempty closed subset of the closed disk in \(\mathbb C\) with center~\(0\) and radius \(\lVert A\rVert\).
\item
State the “Function Representation Theorem” involving an abelian \(C^*\)-algebra \(\mathcal U\) in terms of \(C(\mathcal P(\mathcal U))\), where \(\mathcal P(\mathcal U)\) is the set of all pure states of \(\mathcal U\). Sketch the proof of this theorem.
\item
State and sketch the proof of the theorem involving \(\mathcal A\cong C(X)\), where \(\mathcal A\) is an abelian von Neumann algebra and~\(X\) is an extremely disconnected compact Hausdorff space.
\item
Give the main elements of the Gelfand–Naimark–Segal construction. State the theorem involving a state~\(\rho\) of a \(C^*\)-algebra \(\mathcal U\) and a cyclic representation \(\pi_\rho\) of \(\mathcal U\) on a Hilbert space \(H_\rho\). List the main steps in the proof.
\item
Let \(\mathcal U\) be a \(C^*\)-algebra. State and prove equivalent conditions for \(A\in\mathcal U^+\).
\item
Let~\(H\) be a self-adjoint operator on a Hilbert space. Define the Cayley transform \(U(H)\) and show that \(H\mapsto U(H)\) on the self-adjoint operators is strong-operator continuous.
\item
Let \(\mathcal A\) be the algebra of multiplications by elements of \(L_\infty(S,\mathcal S,m)\) on \(L_2=L_2(S,\mathcal S,m)\). If~\(T\) is a bounded operator on \(L_2\) commuting with \(\mathcal A\), show \(T\in\mathcal A\). Assume \(m(S)<\infty\).
\end{examproblems}
\end{document}
